% =====================================================================
%  front.tex — 封面、译者凡例、摘要
% =====================================================================

% ------------------------- 封面 -------------------------
\begin{titlepage}
\centering
\vspace*{1.2cm}

{\large\bfseries 数译\;·\;数学文献翻译丛书}\\[0.35em]
{\small 第 SY05-25 种}\\[0.6em]
{\color{black!55}\rule{0.62\linewidth}{0.4pt}}

\vspace{2.1cm}

{\Huge\bfseries\linespread{1.15}\selectfont 镜像对称、Hitchin 方程\\[0.5em] 与朗兰兹对偶}\\[1.1em]
{\large\itshape Mirror Symmetry, Hitchin's Equations,\\[0.2em] and Langlands Duality}\\[1.4em]
{\large Edward Witten（爱德华·威滕）}\\[0.4em]
{\small School of Natural Sciences, Institute for Advanced Study, Princeton NJ 08540}

\vspace{1.2cm}
{\color{black!55}\rule{0.62\linewidth}{0.4pt}}

\vfill

\begin{minipage}{0.92\linewidth}
\small\centering
底本：arXiv:0802.0999v1 [math.RT]，2008 年 2 月 7 日，正文 15 页，6 节\\
原文日期：2007 年 2 月\\[0.9em]
中文译注本：语义级忠实翻译 \;+\; 逐符号公式校核 \;+\; nLab 交叉核对的学术译注\\
正文一字未增；一切展开、名词解释与框架梳理，\\[0.2em]
均见脚注「译注」与篇末「译者后记」
\end{minipage}

\vspace{1.4cm}
{\small 数译 \; 译注 \; v1.0 \; 2026 年 11 月}
\vspace*{0.6cm}
\end{titlepage}

% ------------------------- 凡例 -------------------------
\section*{译者凡例}
\addcontentsline{toc}{section}{译者凡例}
\markboth{译者凡例}{译者凡例}

\begin{enumerate}[leftmargin=1.6em,itemsep=0.45em,topsep=0.2em]

\item \textbf{底本与对校。}底本为 arXiv:0802.0999v1 的 PDF（15 页）。翻译时逐段对照原文，公式逐符号与页面图核对；原文的节号、公式号、脚注位置、参考文献条目一律保留原位，公式编号沿用原文的 (2.1)–(5.1)。对校记录见工程文件 \texttt{记录/校核记录.md}。

\item \textbf{翻译原则。}语义级忠实：不改写、不删节、不合并段落、不调整论证次序。作者的第一人称复数（we）、判断的强弱（\textit{it turns out}、\textit{we expect}、\textit{is believed}、\textit{it seems}）与保留态度，一律照译；遇到原文只给结论而不给论证的地方，译文同样只给结论，展开全部放在译注里。

\item \textbf{脚注体例。}全书脚注统一连续编号，分两类：
  \begin{itemize}[leftmargin=1.2em,itemsep=0.1em,topsep=0.1em]
  \item【原注】——作者本人的脚注，原文共 5 条，按其在正文中出现的顺序编号；
  \item【译注 $x.y$】——译者所加，编号中的 $x$ 为所在节号（如【译注 3.2】即第 3 节第 2 条）。
  \end{itemize}
  每条译注末尾以〔〕注明依据，分三级：\textbf{A} = 原文自身或原文所引原始文献；\textbf{B} = nLab 词条（社区维护的百科，均与 A 级交叉核对，并注明访问日期）；\textbf{I} = 本丛书内部文档。凡属译者的推断而无外部依据者，标注「译者按」并说明推断理由，不与上述三级混同。

\item \textbf{术语。}同一英文术语全书只用一个译名，首次出现时在括号内给出原文，其后只用中文；英文术语与中文之间留半角空格。全书术语对照见附录 A。

\item \textbf{符号。}完全沿用原文记号。需注意者三点：复值联络在原文中排粗体，本译本作 $\bm{A}=A+i\phi$，以与实联络 $A$ 区别；Hodge 星算子原文作 $\star$，本译本同；群名 $\SU$、$\PSU$、$\mathrm{E}_8$ 等原文排斜体，本译本统一排正体，属排印统一，不涉内容。

\item \textbf{原文的排印瑕疵。}底本是预印本工作稿，个别词有拼写错误（如 \textit{arbtirarily}、\textit{holomophic}），文末邮箱行把合作者地址与作者地址并列，第 11 页脚注中 \textit{orbifold} 一词的用法亦从上下文才可确定。本译本按语义翻译，不代作者改正；凡影响理解的，另出译注说明，而不改动正文。

\item \textbf{本书的边界。}正文是译文，不是综述：没有增句、没有删句，也没有把作者的猜测写成定理。译者的一切补充——概念解释、公式展开、文献定位、以及篇末对全文框架与推导链条的梳理——都明确标为译者所加。读者若要与原文比对，只需忽略全部【译注】与「译者后记」，剩下的就是原文。

\end{enumerate}

\clearpage

% ------------------------- 摘要 -------------------------
\section*{摘要}
\addcontentsline{toc}{section}{摘要}
\markboth{摘要}{摘要}

\noindent
几何 Langlands 对偶可以从镜像对称的若干陈述来理解，而这些陈述可以对一个定向二维流形 $C$ 用纯拓扑的方式表述。但若不曾在 $C$ 上取定复结构并使用 Hitchin 方程，要理解这些陈述是极其困难的。本文概述两方面的基本陈述：「非分歧」情形，即 $C$ 为紧致无边的定向二维流形；以及「分歧」情形，即允许有穿刺点。我们还略作提示，说明为什么更精确的描述需要以四维规范场论为出发点。

\vspace{0.9em}
\noindent{\small\itshape 关键词：}几何 Langlands 对偶；镜像对称；Hitchin 方程；Higgs 丛；超凯勒几何；分歧；四维规范场论
